\section*{Capítulo \ref{ch:b3:exploration-seismology} --- Sismología de exploración}

\begin{solution}{exo:b3:exploration-seismology:1}
$Z = \rho v$: creta $2000\times2500 = \qty{5.00e6}{kg.m^{-2}.s^{-1}}$, lutita
$2300\times2300 = \qty{5.29e6}{kg.m^{-2}.s^{-1}}$, caliza $2600\times4500 =
\qty{11.70e6}{kg.m^{-2}.s^{-1}}$. Creta--lutita: $R = (5.29 - 5.00)/(5.29 + 5.00)
= 0.29/10.29 = 0.0282$; lutita--caliza: $R = (11.70 - 5.29)/(11.70 + 5.29) =
6.41/16.99 = 0.377$. La energía reflejada es $R^2 = 0.142$: el techo de la
caliza devuelve \qty{14.2}{\percent} de la energía incidente; la interfase creta--lutita
devuelve menos del \qty{0.1}{\percent}.
\end{solution}

\begin{solution}{exo:b3:exploration-seismology:2}
$t_0 = 2\times800/2000 = \qty{0.800}{s}$; $t = \sqrt{0.800^2 + (1000/2000)^2} =
\sqrt{0.89} = \qty{0.9434}{s}$; el sobretiempo es \qty{0.1434}{s}. La
aproximación da $1000^2/(2\times2000^2\times0.800) = \qty{0.1563}{s}$, demasiado
grande en un \qty{9}{\percent}: con una separación mayor que la profundidad, $x \ll vt_0$ ya no
se cumple.
\end{solution}

\begin{solution}{exo:b3:exploration-seismology:3}
Los puntos medios se desplazan la mitad del intervalo entre receptores: el intervalo entre puntos medios comunes es
$12.5/2 = \qty{6.25}{m}$. La cobertura es $240\times12.5/(2\times25) = 60$. Una cobertura de
120 requiere un disparo cada $240\times12.5/(2\times120) = \qty{12.5}{m}$.
\end{solution}

\begin{solution}{exo:b3:exploration-seismology:4}
$\lambda = v/f$: \qty{100}{m} en la arenisca, $\lambda/4 = \qty{25}{m}$;
\qty{150}{m} en la caliza, $\lambda/4 = \qty{37.5}{m}$. Con
$\lambda = 2800/30 = \qty{93.3}{m}$ y $z = \qty{1000}{m}$,
$r_F \approx \sqrt{\lambda z/2} = \sqrt{46\,667} = \qty{216}{m}$ (\qty{217}{m}
con el término $\lambda^2/16$).
\end{solution}

\begin{solution}{exo:b3:exploration-seismology:5}
\omterm{thm:b3:exploration-seismology:dix}{Ecuación de Dix}: $v^2 = (2800^2\times0.700 - 2400^2\times0.500)/0.200 = (5\,488\,000 -
2\,880\,000)/0.200 = \qty{1.304e7}{m^2.s^{-2}}$, luego $v = \qty{3611}{m.s^{-1}}$.
La capa se atraviesa en \qty{0.100}{s} en un sentido: $h = 3611\times0.100 =
\qty{361}{m}$.
\end{solution}

\begin{solution}{exo:b3:exploration-seismology:6}
(a) Sea $d$ la distancia desde el punto medio M al reflector, de modo que
$t_0 = 2d/v$. El rayo reflejado desde la fuente S hasta el receptor G tiene la misma longitud que
la línea recta desde la imagen S$'$ de S en el reflector hasta G. Se sitúa M en
el origen, con la línea según el buzamiento; S$'$ está a la distancia $2d_S$ de S a lo largo
de la normal, con $d_S = d - (x/2)\sin\varphi$. El cuadrado de S$'$G es
$x^2 + 4d_S^2 + 4d_Sx\sin\varphi = x^2 + 4d^2 - x^2\sin^2\varphi$, todos los términos
en $dx$ se cancelan, luego $v^2t^2 = 4d^2 + x^2\cos^2\varphi$ y
$t^2 = t_0^2 + x^2\cos^2\varphi/v^2$. (b) El conjunto de trazas sigue una hipérbola de velocidad
$V = v/\cos\varphi = 2000/\cos\ang{12} = \qty{2045}{m.s^{-1}}$: un buzamiento hace que la
\omterm{def:b3:exploration-seismology:velocities}{velocidad de apilamiento} resulte demasiado alta, aquí en un \qty{2.2}{\percent}.
\end{solution}

\begin{solution}{exo:b3:exploration-seismology:7}
Por $\sqrt{60} = 7.75$. Duplicar la relación señal-ruido requiere cuatro veces tantas
trazas: una cobertura de 240.
\end{solution}

\begin{solution}{exo:b3:exploration-seismology:8}
La profundidad representada $vt/2$ cambia en $\qty{1.5}{km.s^{-1}}\times\qty{0.40}{s}
= \qty{0.60}{km}$ por kilómetro: $\tan\theta_a = 0.60$, $\theta_a = \ang{31.0}$.
Por el teorema de separación nula $\sin\theta = 0.60$: $\theta = \ang{36.9}$.
\end{solution}

\begin{solution}{exo:b3:exploration-seismology:9}
Cada espiga coloca una copia de $w$ centrada en ella: $0.3\,w$ en las muestras 0--2 da
$(-0.15, 0.3, -0.15)$, y $-0.2\,w$ en las muestras 2--4 da $(0.1, -0.2, 0.1)$.
Su suma es $(0, -0.15, 0.3, -0.05, -0.2, 0.1, 0)$, escrita desde una muestra
antes de la serie hasta otra después de ella. Las reflexiones alcanzan su máximo en las muestras 1 y 3, con
sus signos verdaderos; se solapan en la muestra 2, donde $-0.15 + 0.1 = -0.05$: una
ligera interferencia, pero dos máximos distintos.
\end{solution}

\begin{solution}{exo:b3:exploration-seismology:10}
$\arcsin(1800/2900) = \ang{38.37}$ y $\arcsin(1800/4400) = \ang{24.15}$, luego
$i_c = (38.37 + 24.15)/2 = \ang{31.26}$ y $\theta = (38.37 - 24.15)/2 =
\ang{7.11}$; $v_2 = 1800/\sin\ang{31.26} = \qty{3469}{m.s^{-1}}$. La media armónica
$2\times2900\times4400/(2900 + 4400) = \qty{3496}{m.s^{-1}}$ es
un \qty{0.8}{\percent} superior; multiplicada por $\cos\theta = 0.9923$ da
\qty{3469}{m.s^{-1}}, la aproximación habitual de buzamiento pequeño
$v_2 \approx \cos\theta\,2v_dv_u/(v_d + v_u)$.
\end{solution}

\begin{solution}{exo:b3:exploration-seismology:11}
(a) El compuesto es $R[w(t) - w(t - \Delta t)] \approx R\,\Delta t\,w'(t)$ por
el desarrollo de Taylor de primer orden. (b) Su amplitud es proporcional a
$\Delta t$, luego un estrato más delgado que el \omterm{def:b3:exploration-seismology:resolution}{espesor de sintonía} aún devuelve una señal,
con la forma de $w'$, pero su techo y su base no pueden distinguirse como dos eventos: se
detecta, pero no se resuelve. Si se conocen $R$ y la ondícula, por ejemplo a partir de un pozo,
la amplitud da $\Delta t$ y, con ello, el espesor. (c) La amplitud es
proporcional al espesor: la razón es $5/10 = 1/2$.
\end{solution}

\begin{solution}{exo:b3:exploration-seismology:12}
(a) $t(x) = (2/v)\sqrt{z^2 + x^2}$: $t(0) = 2\times600/3000 = \qty{0.400}{s}$ y
$t(300) = (2/3000)\sqrt{600^2 + 300^2} = \qty{0.4472}{s}$. (b) $t^2 = t_0^2 +
x^2/(v/2)^2$, la forma del \omterm{thm:b3:exploration-seismology:nmo}{sobretiempo normal} con la velocidad $v/2$. Una reflexión con el mismo tiempo en el vértice
en una sección de separación nula no mostraría ningún sobretiempo; comparado con una
hipérbola de \omterm{thm:b3:exploration-seismology:nmo}{sobretiempo normal} de igual $t_0$ en un conjunto de trazas, el sobretiempo de la difracción
$x^2/(2(v/2)^2t_0)$ es cuatro veces mayor, porque el tiempo depende de la
distancia al difractor dos veces, de ida y de vuelta. (c) La migración suma la
sección a lo largo de esta hipérbola y sitúa la energía en su vértice: la difracción
se colapsa en un punto en $x = 0$, $z = \qty{600}{m}$.
\end{solution}

\begin{solution}{pb:b3:exploration-seismology:1}
\textbf{1.} Desde O, $400/(486.1 - 425.2) = 400/0.0609 = \qty{6568}{m.s^{-1}}$;
desde E, $400/(486.1 - 372.4) = 400/0.1137 = \qty{3518}{m.s^{-1}}$. La onda de cabeza
desde O parece rápida: O es el extremo más bajo según el buzamiento, y sus disparos viajan hacia arriba del buzamiento; el
refractor asciende de O a E.

\textbf{2.} $\arcsin(2400/3518) = \ang{43.02}$, $\arcsin(2400/6568) =
\ang{21.43}$: $i_c = \ang{32.22}$, $\theta = \ang{10.79}$.

\textbf{3.} $v_2 = 2400/\sin\ang{32.22} = \qty{4501}{m.s^{-1}}$: dentro del
intervalo de la caliza, la cima de la caliza jurásica.

\textbf{4.} Las intercepciones son $t - x/v_{\text{app}}$ a \qty{1200}{m}:
$t_{i,O} = 0.4861 - 1200/6568 = \qty{0.3034}{s}$ y $t_{i,E} = 0.4861 -
1200/3518 = \qty{0.1450}{s}$. Con $\cos i_c = 0.8460$, $h = v_1t_i/(2\cos
i_c)$: $h_O = 2400\times0.3034/1.692 = \qty{430.4}{m}$ y $h_E =
\qty{205.7}{m}$.

\textbf{5.} Dividiendo entre $\cos\theta = 0.9823$: \qty{438.1}{m} bajo O y
\qty{209.4}{m} bajo E. Entonces $(438.1 - 209.4)/1200 = 0.1906 =
\tan\ang{10.79}$, el buzamiento de la pregunta 2.

\textbf{6.} La lutita no devuelve onda de cabeza: es invisible, y la prospección
la cuenta junto con la creta. Aquí $v_1 = \qty{2.40}{km.s^{-1}}$ se dio como
media de las dos; si para todo el recubrimiento se hubiera usado la velocidad de la
onda directa de la creta, \qty{2.50}{km.s^{-1}}, todas las profundidades saldrían demasiado grandes.

\textbf{7.} $V = 600/\sqrt{0.5437^2 - 0.4878^2} = 600/0.2401 =
\qty{2499}{m.s^{-1}}$ y $V = 600/\sqrt{0.6240^2 - 0.5748^2} =
\qty{2471}{m.s^{-1}}$. La segunda promedia la creta con la lutita, más lenta,
que queda debajo: la media desciende.

\textbf{8.} \qty{2811}{m.s^{-1}} y \qty{2829}{m.s^{-1}}.

\textbf{9.} Capa 1: \qty{2499}{m.s^{-1}}. Capa 2: $v^2 = (2471^2\times0.5748
- 2499^2\times0.4878)/0.0870$, $v = \qty{2308}{m.s^{-1}}$. Capa 3:
$v^2 = (2811^2\times0.6548 - 2471^2\times0.5748)/0.0800$,
$v = \qty{4561}{m.s^{-1}}$. Capa 4: $v^2 = (2829^2\times0.7348 -
2811^2\times0.6548)/0.0800$, $v = \qty{2972}{m.s^{-1}}$.

\textbf{10.} Creta (\qty{2499}{m.s^{-1}}), lutita (\qty{2308}{m.s^{-1}}),
caliza (\qty{4561}{m.s^{-1}}), arenisca (\qty{2972}{m.s^{-1}}), cada una dentro
de su intervalo publicado. Las reflexiones son la interfase creta--lutita, la
cima de la caliza, la interfase caliza--arenisca y la discordancia en
la cima del basamento de esquistos.

\textbf{11.} La creta y la lutita tienen casi la misma impedancia: $R = 0.028$
(ejercicio 1), frente a 0.38 en la cima de la caliza.

\textbf{12.} La hipérbola de velocidad $V_{\text{rms}}$ solo se cumple hasta el
segundo orden en el desplazamiento; a desplazamientos de \qty{2}{km}, el doble
de la profundidad del reflector más profundo, los tiempos se apartan de ella y
una hipérbola ajustada a todo el tendido daría velocidades erróneas, y los
intervalos de Dix, aún peores. Los desplazamientos lejanos estiran también las ondículas al corregirlas.

\textbf{13.} Tiempos de ida y vuelta en las capas: \qty{0.4878}{s}, \qty{0.0870}{s},
\qty{0.0800}{s}, \qty{0.0800}{s}. Espesores $v\,\Delta t/2$: \qty{609.5}{m},
\qty{100.4}{m}, \qty{182.4}{m}, \qty{118.9}{m}.

\textbf{14.} Profundidades bajo el suelo: \qty{609.5}{m}, \qty{709.9}{m},
\qty{892.3}{m}, \qty{1011.2}{m}. La lutita, la caliza y la arenisca salen
dentro de \qty{0.4}{m}, \qty{2.4}{m} y \qty{1.1}{m} de sus espesores, creciendo
el error en la caliza, rápida y delgada, donde un pequeño error de tiempo
cuesta más.

\textbf{15.} $\bar v = 2\times1011.2/0.7348 = \qty{2752}{m.s^{-1}}$. La
\omterm{def:b3:exploration-seismology:velocities}{velocidad cuadrática media} pondera cada capa por $v^2$ y nunca es inferior a la
media ponderada por el tiempo: $2829\times0.7348/2 = \qty{1039}{m}$, \qty{28}{m} demasiado profunda.

\textbf{16.} $\lambda = 2308/30 = \qty{76.9}{m}$, $\lambda/4 = \qty{19.2}{m}$;
la estimación de Ricker es $2308/(5.2\times30) = \qty{14.8}{m}$. La lutita tiene
\qty{100}{m} de espesor, cinco veces su \omterm{def:b3:exploration-seismology:resolution}{espesor de sintonía}: su cima y su base
se resuelven. Sintonizaría, según la regla del cuarto de longitud de onda, a $f = v/(4b) =
2308/400 = \qty{5.77}{Hz}$.

\textbf{17.} $\lambda = 2829/30 = \qty{94.3}{m}$, $z = \qty{1011}{m}$:
$r_F \approx \sqrt{94.3\times1011/2} = \qty{218}{m}$ (\qty{220}{m} con el
término $\lambda^2/16$).

\textbf{18.} $t(x)^2 = t_0^2 + (2x/v)^2$, así que $v = 2x/\sqrt{t^2 - t_0^2} =
400/\sqrt{0.5732^2 - 0.550^2} = 400/0.1614 = \qty{2478}{m.s^{-1}}$.

\textbf{19.} $z = vt_0/2 = 2478\times0.550/2 = \qty{681}{m}$.

\textbf{20.} La profundidad representada $vt/2$ cambia $1.24\times0.167 = 0.207$ por
unidad de longitud: $\theta_a = \arctan0.207 = \ang{11.70}$. El buzamiento verdadero
cumple $\sin\theta = 0.207$: $\theta = \ang{11.95}$, el buzamiento del flanco.

\textbf{21.} $d = 2480\times0.655/2 = \qty{812}{m}$; la migración desplaza la
reflexión $d\sin\theta = 812\times0.207 = \qty{168}{m}$ hacia arriba del buzamiento, hacia la
cresta del anticlinal, y hacia arriba desde \qty{812}{m} hasta $d\cos\theta =
\qty{795}{m}$.

\textbf{22.} $2500\times0.160/2 = \qty{200}{m}$.

\textbf{23.} De las preguntas 13 y 14 se deduce que la discordancia del
basamento se halla \textbf{\qty[text-series-to-math=true]{1011}{m}} bajo el suelo en el \omterm{def:b3:exploration-seismology:cmp}{punto medio común} S.
\end{solution}
